\section{Motivation}
% Databases form the backbone of nearly every modern software system, supporting essential operations such as data storage, retrieval, analysis, and application logic. As organizations increasingly rely on data-driven decision making, the need for efficient, scalable, and accessible database management systems has become more critical than ever. Over the years, a wide variety of database technologies have emerged, ranging from traditional relational databases (SQL) such as PostgreSQL and MySQL to non-relational databases (NoSQL) such as MongoDB and Cassandra. Each technology introduces unique features, use cases, and architectural principles, enabling developers to choose the most suitable database model for their applications. 

%  However, this diversity also brings substantial challenges. Developers and students must work with multiple tools, interfaces, and frameworks—often switching between environments for schema design, data manipulation, cloud hosting, and query execution. This fragmented workflow increases complexity, slows down productivity, and creates a steep learning curve, especially for individuals who are new to database development. As a result, many teams struggle with consistency, collaboration, and maintaining clear understanding of database structures across different projects. 

% With the rise of cloud-based development practices, the industry has seen a shift toward managed services such as AWS RDS, Supabase, Firebase, and MongoDB Atlas. While these platforms simplify certain operational aspects, they primarily focus on hosting and infrastructure rather than providing a fully integrated development and management environment. Developers still need several external tools for schema visualization, data entry, query editing, version control, and project automation. This creates inefficiencies that slow down development cycles and make database management more error-prone. 

% At the same time, advancements in artificial intelligence have opened new opportunities for improving developer productivity. AI-powered systems can assist in generating schemas, understanding complex structures, writing optimized queries, and interacting with databases through natural language. Techniques such as Retrieval-Augmented Generation (RAG) allow systems to interpret user intent more accurately and provide context-aware answers based on the actual structure and content of a database. Despite these capabilities, existing database tools rarely integrate AI in a meaningful or accessible way. 

% In response to these challenges, there is a growing need for a unified, intelligent Database-as-a-Service (DBaaS) platform that combines the management of SQL and NoSQL databases within a single, intuitive interface. Such a platform would streamline workflow, eliminate tool fragmentation, accelerate development, lower the learning curve for beginners, and empower developers to focus on building applications rather than managing database infrastructure. By incorporating AI-powered assistance, automated schema generation, cloud storage, versioning, and query guidance, the platform can significantly enhance both efficiency and accessibility. 

% This project aims to address these needs by developing an integrated, cloud-based DBaaS system that supports PostgreSQL and MongoDB, provides visual schema design, enables seamless data manipulation, and incorporates AI through Retrieval-Augmented Generation. The platform is designed to serve as a comprehensive and user-friendly solution that simplifies database development for both beginners and experienced engineers.

Databases are essential to modern software systems, supporting data storage, retrieval, analysis, and application logic. The increasing reliance on data-driven decision-making has elevated the need for efficient, scalable, and accessible database management systems, including relational databases (SQL) like PostgreSQL and MySQL, and non-relational databases (NoSQL) such as MongoDB and Cassandra.

Database development, however, is often complex due to fragmented workflows involving multiple tools for schema design, data manipulation, query execution, and cloud hosting. While cloud-based services like AWS RDS, Supabase, Firebase, and MongoDB Atlas simplify infrastructure management, they do not provide fully integrated development environments, leaving inefficiencies and steep learning curves for developers.

Advances in artificial intelligence, including Retrieval-Augmented Generation (RAG), offer opportunities for automated schema generation, optimized queries, and natural language interactions, yet AI integration in existing tools remains limited.

This project addresses these challenges by developing a unified, cloud-based Database-as-a-Service (DBaaS) platform supporting PostgreSQL and MongoDB. It combines visual schema design, seamless data manipulation, cloud storage, and AI-powered assistance, providing an intuitive, efficient, and accessible environment for both novice and experienced developers.